\begin{scanpage}{0114}{97}
not succeed in obtaining correct results for a cone with an arbitrary aperture angle $\omega$ $(0\leq\omega\leq\pi/2)$.

The effective scattering area in accordance with (12.18) is determined by the equation
\[\sigma=\pi a|\Sigma|^2,\tag{17.06}\]
where the function $\Sigma$ is connected with the fringing field by the relationship
\[E_x=-H_y=-\frac a2E_{0x}\frac{e^{ikR}}R\Sigma\tag{17.07}\]
and equals
\[\Sigma=\frac1{ka}\operatorname{tg}^2\omega\sin kl\,e^{ikl}+\frac{\frac2n\sin\frac\pi n}{\cos\frac\pi n-\cos\frac{2\omega}n}e^{2ikl}.\tag{17.08}\]
The analogous function in the physical optics approach may be written in accordance with (17.02) in the form
\[\Sigma^0=\frac1{ka}\operatorname{tg}^2\omega\sin kl\,e^{ikl}-\operatorname{tg}\omega\,e^{2ikl}.\tag{17.09}\]
With the deforming of the top part of the cone into a disk $(\omega\to\frac\pi2,\ l\to0)$, Equations (17.08) and (17.09) are transformed, respectively, to the form
\[\left.\begin{aligned}\Sigma&=-ika-\frac1n\operatorname{ctg}\frac\pi n,\\\Sigma^0&=-ika.\end{aligned}\right\}\tag{17.10}\]
Furthermore, it follows from (17.08) and (17.09) that for large values of the parameter $ka$ $(ka\gg\operatorname{tg}^2\omega)$ the functions $\Sigma$ and $\Sigma^0$ may be represented in the form
\[\Sigma=\frac{\frac2n\sin\frac\pi n}{\cos\frac\pi n-\cos\frac{2\omega}n}e^{2ikl},\tag{17.11}\]
\end{scanpage}
